% =====================================================================
%  Lecture 21 -- Network Slicing                           (Persian edition)
% =====================================================================
\documentclass[11pt,a4paper]{article}
\providecommand\ACNroot{../..}
\input{\ACNroot/templates/lecture-fa}

\ACNset{number}{21}
\ACNset{title}{برش‌بندی شبکه}
\ACNset{shorttitle}{برش‌بندی شبکه}
\ACNset{titleother}{Network Slicing}
\ACNset{part}{بخش چهارم -- سامانه و رویه‌های نسل پنجم (درس‌های ۱۴ و ۱۹ تا ۲۲)}
\ACNset{time}{۶ تا ۸ ساعت (۳ ساعت مطالعه، ۳ تا ۵ ساعت تمرین)}
\ACNset{level}{پیشرفته}
\ACNset{prereq}{درس‌های ۲ (ارلانگ \lr{B})، ۱۶، ۱۹ و ۲۰؛ توزیع نرمال}
\ACNset{gen}{5}

\begin{document}
\ACNcover

% ---------------------------------------------------------------------
\section{اهداف یادگیری}
پس از مطالعهٔ این درس باید بتوانید:
\begin{objectives}
  \item برش شبکه را \textbf{تعریف کنید} و \lr{S-NSSAI} و مجموعه‌های \acro{NSSAI} به‌کاررفته در سیگنالینگ را \textbf{رمزگشایی کنید}.
  \item انتخاب برش هنگام ثبت‌نام و برقراری نشست \acro{PDU} ویژهٔ برش را \textbf{دنبال کنید}.
  \item \textbf{توضیح دهید} برش‌ها چگونه در \acro{RAN}، انتقال و هسته محقق می‌شوند و گزینه‌های جداسازی چیستند.
  \item مدیریت برش -- چرخهٔ عمر، الگوها و نقش‌های اپراتور و مستأجر -- را \textbf{توصیف کنید}.
  \item هزینهٔ چندگانه‌سازی آماری جداسازی، بخش‌بندی منابع \acro{RAN} و کنترل پذیرش برش را \textbf{کمّی کنید}.
\end{objectives}

% ---------------------------------------------------------------------
\section{دانش پیش‌نیاز}
به معماری و رویه‌های \lr{5GC} (درس‌های ۱۹ و ۲۰)، مفاهیم مجازی‌سازی درس ۱۶ و فرمول ارلانگ \lr{B} درس ۲ نیاز دارید. مدل
چندگانه‌سازی از توزیع نرمال استفاده می‌کند: برای باری گاوسی با میانگین $\mu$ و انحراف معیار $\sigma$، ظرفیتی که با احتمال $\epsilon$
از آن فراتر رفته می‌شود $\mu + z_{1-\epsilon}\sigma$ است ($z_{0.99} = 2.326$).

\paragraph{واژه‌های کلیدی} واژگان این درس همراه با معادل انگلیسی آن‌ها:
\begin{keyterms}
برش شبکه & \lr{Network slice} & شبکهٔ منطقی‌ای که قابلیت‌ها و ویژگی‌های مشخصی فراهم می‌کند. \\
شناسهٔ برش & \lr{S-NSSAI} & اطلاعات کمکی انتخاب یک برش شبکه (\lr{SST + SD}). \\
نوع برش / متمایزکننده & \lr{SST / SD} & نوع برش/خدمت و متمایزکنندهٔ اختیاری برش. \\
مجموعهٔ \lr{S-NSSAI}ها & \lr{NSSAI} & مجموعه‌های پیکربندی‌شده، درخواستی، مجاز و ردشده. \\
تابع انتخاب برش & \lr{NSSF} & \acro{NF}ای که برش‌های مجاز و مجموعهٔ \acro{AMF} را انتخاب می‌کند. \\
احراز هویت ویژهٔ برش & \lr{NSSAA} & احراز هویت ویژهٔ برش توسط یک کارساز \acro{AAA} شخص ثالث. \\
کنترل پذیرش برش & \lr{NSAC} & محدودیت‌های تعداد \acro{UE} یا نشست \acro{PDU} در هر برش. \\
جداسازی & \lr{Isolation} & میزان محافظت یک برش از برش دیگر. \\
مستأجر & \lr{Tenant} & مشتری (مثلاً یک سازمان) که برش برای او اداره می‌شود. \\
الگوی برش & \lr{Slice template} & توصیف استانداردشدهٔ ویژگی‌های برش (مثلاً تأخیر، گذردهی). \\
\end{keyterms}

% ---------------------------------------------------------------------
\section{مقدمه و انگیزه}
سه سناریوی \lr{IMT-2020} (درس ۱۴) خواسته‌های متضادی دارند: یک خدمت \acro{eMBB} گذردهی می‌خواهد و تأخیر را تحمل می‌کند، یک خدمت
\acro{URLLC} تأخیر و قابلیت اطمینان تضمین‌شده می‌خواهد و یک خدمت \acro{mMTC} تعداد عظیمی اتصال ارزان و به‌ندرت فعال. ساختن سه
شبکهٔ جداگانه غیراقتصادی است؛ اجرای همه روی یک شبکه با اشتراک بهترین تلاش، سخت‌گیرترین مشتریان را ناکام می‌گذارد.
\emph{برش‌بندی شبکه} چند شبکهٔ منطقی را روی یک زیرساخت فیزیکی اجرا می‌کند که هر یک قابلیت‌ها، سیاست‌ها و در صورت نیاز کارکردهای
شبکهٔ خود را دارد \cite{ts23501,foukas2017slicing}.

برش‌بندی یک مفهوم تجاری نیز هست: اپراتور می‌تواند برشی را به‌عنوان محصولی با ویژگی‌های تضمین‌شده به یک سازمان، یک نهاد ایمنی
عمومی یا یک صنعت عمودی عرضه کند که از طریق واسط‌های استاندارد مدیریت می‌شود \cite{ts28530}. این کار با فناوری‌های بخش سوم ممکن
می‌شود: مجازی‌سازی نمونه‌های اختصاصی را ارزان می‌سازد، \acro{SDN} انتقال را برنامه‌ریزی می‌کند و \acro{SBA} اجازه می‌دهد
\acro{NF}ها برای هر برش انتخاب شوند.

% ---------------------------------------------------------------------
\section{مفاهیم بنیادی}

\subsection{شناسایی برش‌ها}
هر برش با یک \lr{S-NSSAI} شناسایی می‌شود (\cref{fig:snssai}): یک \emph{نوع برش/خدمت} (\acro{SST}) ۸ بیتی و یک \emph{متمایزکنندهٔ برش}
(\acro{SD}) اختیاری ۲۴ بیتی که چند برش از یک نوع را از هم جدا می‌کند (مثلاً برش‌های \acro{URLLC} دو سازمان) \cite{ts23003}. مقادیر
استانداردشدهٔ \acro{SST} شامل ۱ (\acro{eMBB})، ۲ (\acro{URLLC})، ۳ (\acro{IoT} انبوه) و ۴ (\acro{V2X}) است و مقادیر دیگری در انتشارهای
بعدی افزوده شده‌اند؛ اپراتورها می‌توانند مقادیر خود را نیز تعریف کنند. یک \acro{UE} با مجموعه‌هایی از \lr{S-NSSAI}ها کار می‌کند:
\begin{itemize}
  \item \textbf{\acro{NSSAI} پیکربندی‌شده}: برای هر \acro{PLMN} در \acro{UE} فراهم می‌شود؛
  \item \textbf{\acro{NSSAI} درخواستی}: در \lr{Registration Request} فرستاده می‌شود؛
  \item \textbf{\acro{NSSAI} مجاز}: توسط شبکه برای ناحیهٔ ثبت‌نام برگردانده می‌شود (حداکثر ۸ \lr{S-NSSAI})؛
  \item \textbf{\acro{NSSAI} ردشده}: برش‌های درخواستی ردشده، همراه با علت (مثلاً در این ناحیه در دسترس نیست).
\end{itemize}

\begin{acnfigure}
  \resizebox{\linewidth}{!}{\acnfig{lecture-21/snssai}}
  \caption{ساختار \lr{S-NSSAI} و مجموعه‌های \acro{NSSAI} مبادله‌شده هنگام ثبت‌نام.}
  \label{fig:snssai}
\end{acnfigure}

\subsection{یک برش از چه تشکیل شده است}
یک برش سرتاسری \acro{RAN}، شبکهٔ انتقال و هسته را دربر می‌گیرد (\cref{fig:slices}). برخی کارکردها معمولاً مشترک‌اند (\acro{AMF}،
\acro{NSSF}، \acro{NRF}، \acro{UDM}، گره‌های \acro{RAN}) و برخی ممکن است ویژهٔ برش باشند (\acro{SMF}، \acro{UPF}، \acro{PCF}،
کارکردهای کاربرد). یک \acro{UE} می‌تواند هم‌زمان تا هشت برش را به کار ببرد، اما همه از طریق یک \acro{AMF}؛ هر نشست \acro{PDU}
دقیقاً به یک برش تعلق دارد.

\begin{acnfigure}
  \resizebox{\linewidth}{!}{\acnfig{lecture-21/slices}}
  \caption{سه برش سرتاسری که \acro{RAN} و کارکردهای کنترلی مشترک را به اشتراک می‌گذارند، با \acro{SMF}ها و \acro{UPF}های ویژهٔ برش.}
  \label{fig:slices}
\end{acnfigure}

\subsection{جداسازی}
جداسازی در چند سطح ممکن است: \emph{منطقی} (سیاست‌ها و کیفیت خدمات جداگانه روی \acro{NF}های مشترک)، \emph{مجازی} (نمونه‌های
\acro{NF} اختصاصی روی زیرساخت مشترک) یا \emph{فیزیکی} (سخت‌افزار، طیف یا سایت اختصاصی). جداسازی قوی‌تر یک برش را از بار، خطاها و
حوادث امنیتی دیگران محافظت می‌کند، اما بهرهٔ چندگانه‌سازی آماری اشتراک را از دست می‌دهد (\cref{sec:mux}). سطح مناسب به توافق‌نامهٔ
سطح خدمت بستگی دارد: یک برش ایمنی عمومی ممکن است به \acro{UPF}های اختصاصی و منابع رادیویی رزروشده نیاز داشته باشد، در حالی که یک
برش ویدئوی مصرفی می‌تواند تقریباً همه‌چیز را به اشتراک بگذارد.

% ---------------------------------------------------------------------
\section{توضیحات فنی تفصیلی}

\subsection{انتخاب برش}
هنگام ثبت‌نام، \acro{AMF} \acro{NSSAI} درخواستی را با \lr{S-NSSAI}های مشترک‌شده از \acro{UDM} مقایسه می‌کند و در صورت لزوم از
\acro{NSSF} می‌پرسد که \acro{NSSAI} مجاز برای موقعیت \acro{UE} و احتمالاً مجموعهٔ \acro{AMF} متفاوتی که برای برش‌ها مناسب‌تر است را
برمی‌گرداند (\cref{fig:sel}). آن‌گاه \acro{AMF} ممکن است ثبت‌نام را به یک \acro{AMF} از آن مجموعه هدایت کند. وقتی \acro{UE} برای یک
\lr{S-NSSAI} و \acro{DNN} نشست \acro{PDU} درخواست می‌کند، \acro{AMF} از طریق \acro{NRF} یک \acro{SMF} خدمت‌دهنده به آن برش را کشف
می‌کند و \acro{SMF} یک \acro{UPF} از آن برش برمی‌گزیند. اطلاعات برش به \acro{RAN} نیز داده می‌شود که ممکن است برای تخصیص منابع و
تصمیم‌های تحرک از آن استفاده کند.

\begin{acnfigure}
  \resizebox{\linewidth}{!}{\acnfig{lecture-21/selection}}
  \caption{انتخاب برش هنگام ثبت‌نام و برقراری نشست \acro{PDU} ویژهٔ برش (ساده‌شده).}
  \label{fig:sel}
\end{acnfigure}

\subsection{احراز هویت ویژهٔ برش و کنترل پذیرش}
انتشار ۱۶ \emph{احراز هویت و مجوزدهی ویژهٔ برش شبکه} (\acro{NSSAA}) را افزود: پس از احراز هویت اولیه، دسترسی به یک برش می‌تواند
افزون بر آن به اعتبارنامه‌هایی وابسته باشد که یک کارساز \acro{AAA} مستأجر برش (مثلاً یک سازمان) بررسی می‌کند، تا مستأجر کنترل کند
کدام دستگاه‌ها از برش او استفاده کنند. انتشار ۱۷ \emph{کنترل پذیرش برش شبکه} (\acro{NSAC}) را افزود: شبکه می‌تواند تعداد
\acro{UE}های ثبت‌شده و نشست‌های \acro{PDU} را در هر برش محدود کند و ظرفیتی را که به یک مستأجر فروخته اعمال کند \cite{ts23501}.

\subsection{برش‌بندی در \lr{RAN} و انتقال}
\acro{RAN} مشترک است؛ برش‌ها با سیاست‌های زمان‌بندی محقق می‌شوند: سهم‌های کمینه، بیشینه و اختصاصی از بلوک‌های منبع فیزیکی برای هر برش،
کنترل پذیرش و اولویت آگاه از برش، و احتمالاً بخش‌های پهنای باند یا حامل‌های جداگانه. شبکهٔ انتقال برش‌ها را به شبکه‌های خصوصی
مجازی، سیاست‌های مسیریابی قطعه‌ای یا کانال‌های جداسازی‌شدهٔ سخت با پهنای باند و تأخیر تضمین‌شده نگاشت می‌کند. بدون این‌ها، یک برش
هسته به‌تنهایی هیچ تضمین سرتاسری‌ای نمی‌دهد.

\subsection{مدیریت برش}
مشخصات مدیریتی \acro{3GPP} برش‌ها را به‌صورت اشیای مدیریت‌شده توصیف می‌کنند: یک \emph{برش شبکه} متشکل از \emph{زیربرش‌های شبکه}
(\acro{RAN}، انتقال، هسته)، با چرخهٔ عمر آماده‌سازی، راه‌اندازی، بهره‌برداری و برچیدن \cite{ts28530}. نیازمندی‌های مشتریان با الگوهای
برش بیان می‌شوند -- برای نمونه الگوی کلی برش شبکهٔ \acro{GSMA} ویژگی‌هایی مانند گذردهی، تأخیر، دسترس‌پذیری، تعداد دستگاه‌ها و
جداسازی را فهرست می‌کند -- که هماهنگ‌سازی اپراتور آن‌ها را با استفاده از \acro{MANO} در \acro{NFV} (درس ۱۶) و کنترل‌گرهای \acro{SDN}
(درس ۱۵) به نمونه‌های \acro{NF}، سیاست‌ها و پیکربندی \acro{RAN} و انتقال ترجمه می‌کند.

% ---------------------------------------------------------------------
\section{مدل‌های ریاضی}

\subsection{هزینهٔ جداسازی}\label{sec:mux}
\paragraph{مدل} برش $i$ باری گاوسی با میانگین $\mu_i$ و انحراف معیار $\sigma_i$ دارد که از دیگران مستقل است. برای نگه‌داشتن احتمال
اضافه‌بار زیر $\epsilon$، ظرفیت اختصاصی هر برش و ظرفیت مشترک همهٔ برش‌ها عبارت‌اند از
\begin{equation}\label{eq:dedicated}
  C_i = \mu_i + z\,\sigma_i, \qquad
  C_{\mathrm{shared}} = \sum_i \mu_i + z\sqrt{\sum_i \sigma_i^2},
\end{equation}
با $z = z_{1-\epsilon}$. \emph{بهرهٔ چندگانه‌سازی} اشتراک
\begin{equation}\label{eq:gain}
  G = \frac{\sum_i C_i}{C_{\mathrm{shared}}} \ge 1
\end{equation}
است.
\paragraph{مثال} سه برش با $\mu = 100$ و $\sigma = 30$ (در هر واحد ظرفیتی) و $\epsilon = 1\,\%$ ($z = 2.326$): اختصاصی $3\times169.8 = 509.3$،
مشترک $300 + 2.326\sqrt3\times30 = 420.9$؛ $G = 1.21$ -- جداسازی ۲۱٪ ظرفیت بیشتر هزینه دارد (\cref{fig:mux}).
\paragraph{تفسیر} بهره با تعداد برش‌ها و با انفجاری‌بودن $\sigma/\mu$ رشد می‌کند و برای بارهای ثابت ناپدید می‌شود. بنابراین جداسازی
دقیقاً برای خدمات انفجاری گران است -- استدلال اقتصادی برای اشتراک آنچه می‌توان به اشتراک گذاشت و جداسازی فقط آنچه توافق‌نامهٔ سطح
خدمت لازم می‌داند. بارهای همبسته (مثلاً اوج همهٔ برش‌ها در یک ساعت) بهره را کاهش می‌دهند.

\begin{acnfigure}
  \acnfig{lecture-21/mux-gain}
  \caption{بهرهٔ چندگانه‌سازی یک مخزن مشترک در برابر ظرفیت اختصاصی برش‌ها (\cref{eq:gain}، $\epsilon = 1\,\%$).}
  \label{fig:mux}
\end{acnfigure}

\subsection{بخش‌بندی منابع \lr{RAN}}\label{sec:prb}
\paragraph{مدل} یک سلول $N$ بلوک منبع دارد؛ برش $i$ سهم کمینهٔ $\rho_i$ را تضمین‌شده دارد و با یک طرح مدولاسیون و کدگذاری مقاوم، نرخ
$r_i$ برای هر بلوک منبع. گذردهی تضمین‌شدهٔ آن
\begin{equation}\label{eq:prb}
  R_i^{\min} = \lfloor\rho_i N\rfloor\, r_i
\end{equation}
است، در حالی که منابع تضمین‌شدهٔ استفاده‌نشده می‌توانند به برش‌های دیگر قرض داده شوند مگر آنکه \emph{اختصاصی} باشند.
\paragraph{مثال} $N = 273$ (۱۰۰ مگاهرتز با ۳۰ کیلوهرتز)، کمینهٔ \acro{URLLC} ۲۰٪ (۵۴ بلوک منبع) و $r = 0.5$ مگابیت بر ثانیه برای هر بلوک:
$R^{\min} = 27$ مگابیت بر ثانیه، کافی برای بسیاری از جریان‌های کنترلی، در حالی که بقیهٔ سلول هنگام بیکاری \acro{URLLC} در اختیار
\acro{eMBB} است.

\subsection{کنترل پذیرش برش}\label{sec:nsac}
\paragraph{مدل} اگر برش یک مستأجر به $c$ نشست \acro{PDU} هم‌زمان محدود باشد و نشست‌ها به‌صورت جریان پواسون با بار ارائه‌شدهٔ $A$
ارلانگ برسند، احتمال رد یک نشست جدید فرمول ارلانگ \lr{B} درس ۲ است:
\begin{equation}\label{eq:block}
  B(A, c) = \frac{A^{c}/c!}{\sum_{k=0}^{c} A^{k}/k!} .
\end{equation}
\paragraph{مثال} $A = 40$ ارلانگ: با $c = 50$، $B = 1.9\,\%$؛ با $c = 60$، $B = 0.07\,\%$. سهمیهٔ $c = 53$ کوچک‌ترین مقداری است که $B \le 1\,\%$
می‌دهد. سهمیه بخشی از قرارداد است: از برش‌های دیگر محافظت می‌کند و تعیین می‌کند مستأجر برای چه پول می‌پردازد.

% ---------------------------------------------------------------------
\section{معماری شبکه و نمودارها}
\Cref{fig:slices} یک الگوی استقرار رایج را نشان می‌دهد: \acro{RAN} مشترک با سیاست‌های زمان‌بندی برای هر برش، \acro{VPN}های انتقال برای
هر برش، \acro{AMF} و دادهٔ مشترک مشترک، و \acro{SMF} و \acro{UPF} برای هر برش -- \acro{UPF} برش \acro{URLLC} در لبه و \acro{UPF}های برش
\acro{eMBB} در مرکز. یک برش سازمانی ممکن است افزون بر آن یک \acro{UPF} در محل سازمان (درس ۱۹) و \acro{NSSAA} به سوی کارساز \acro{AAA}
سازمان داشته باشد. سامانهٔ مدیریت نگاشت هر برش به زیربرش‌ها و منابعش را نگه می‌دارد.

% ---------------------------------------------------------------------
\section{رویه‌های گام‌به‌گام}

\begin{procedurebox}{ثبت‌نام با انتخاب برش}\label{proc:sel}
\begin{steps}
  \item \acro{UE} یک \lr{Registration Request} با \acro{NSSAI} درخواستی (از \acro{NSSAI} پیکربندی‌شدهٔ خود) می‌فرستد.
  \item \acro{AMF} \lr{S-NSSAI}های مشترک‌شده را از \acro{UDM} بازیابی می‌کند.
  \item اگر به‌تنهایی نتواند تصمیم بگیرد، \acro{AMF} با \acro{NSSAI} درخواستی و مشترک‌شده و \acro{TAI} \acro{UE} از \acro{NSSF} پرس‌وجو
        می‌کند؛ \acro{NSSF} \acro{NSSAI} مجاز و احتمالاً یک مجموعهٔ \acro{AMF} هدف را برمی‌گرداند.
  \item اگر مجموعهٔ \acro{AMF} دیگری لازم باشد، ثبت‌نام هدایت می‌شود؛ برش‌هایی که به \acro{NSSAA} نیاز دارند با کارساز \acro{AAA} مستأجر
        مجوز می‌گیرند.
  \item \lr{Registration Accept} \acro{NSSAI} مجاز و ردشده را برمی‌گرداند؛ \acro{UE} در درخواست‌های نشست \acro{PDU} فقط از
        \lr{S-NSSAI}های مجاز استفاده می‌کند.
\end{steps}
\end{procedurebox}

\begin{procedurebox}{راه‌اندازی یک برش سازمانی}\label{proc:commission}
\begin{steps}
  \item نیازمندی‌های مستأجر در یک الگوی برش ثبت می‌شوند (گذردهی، تأخیر، دستگاه‌ها، ناحیه، جداسازی، دسترس‌پذیری).
  \item سامانهٔ مدیریت اپراتور برش را به زیربرش‌های \acro{RAN}، انتقال و هسته تجزیه و امکان‌پذیری را بررسی می‌کند.
  \item هماهنگ‌سازی \acro{NFV} \lr{SMF/UPF} ویژهٔ برش (مثلاً یک \acro{UPF} درون‌محلی) را نمونه‌سازی می‌کند؛ کنترل‌گرهای \acro{SDN}
        مسیرهای انتقال را برقرار می‌کنند؛ سیاست‌های \acro{RAN} سهم منابع را در سلول‌های تحت پوشش رزرو می‌کنند.
  \item اشتراک دستگاه‌های مستأجر با \lr{S-NSSAI} جدید به‌روز می‌شود؛ پیکربندی \acro{NSSF} و سهمیه‌های \acro{NSAC} به‌روز می‌شوند.
  \item برش فعال و کارایی آن در برابر توافق‌نامهٔ سطح خدمت پایش می‌شود.
\end{steps}
\end{procedurebox}

% ---------------------------------------------------------------------
\section{مثال‌های حل‌شده}

\begin{examplebox}{رمزگشایی \lr{S-NSSAI}ها}\label{ex:snssai}
\textbf{مسئله.} \acro{NSSAI} مجاز یک \acro{UE} شامل (\lr{SST 1})، (\lr{SST 2}، \lr{SD \texttt{0x00A1B2}}) و (\lr{SST 2}، \lr{SD \texttt{0x00C3D4}}) است.
آن را تفسیر کنید.

\textbf{حل.} یک برش پیش‌فرض \acro{eMBB} بدون متمایزکننده و دو برش \acro{URLLC} که با \acro{SD}هایشان متمایز می‌شوند -- مثلاً دو مشتری سازمانی یا
صنعتی مختلف. \acro{UE} می‌تواند هم‌زمان در هر سه برش از طریق همان \acro{AMF} نشست \acro{PDU} باز کند.
\end{examplebox}

\begin{examplebox}{هزینهٔ جداسازی}\label{ex:mux}
\textbf{مسئله.} با \cref{eq:dedicated}، اپراتور با اشتراک ظرفیت \acro{UPF} میان سه برش \cref{sec:mux} به‌جای اختصاص آن چه مقدار ظرفیت صرفه‌جویی
می‌کند و مستأجر \acro{URLLC} چه خطری می‌پذیرد؟

\textbf{حل.} $509.3 - 420.9 = 88.4$ واحد (۱۷٪ کل اختصاصی). آن‌گاه برش \acro{URLLC} در اوج‌های مشترک با دیگران رقابت می‌کند؛ اگر توافق‌نامهٔ آن
کارایی تضمین‌شده مستقل از ترافیک دیگران بخواهد، سهم آن باید رزرو (جداسازی منطقی با اولویت) یا اختصاصی (جداسازی مجازی) شود.
\end{examplebox}

\begin{examplebox}{اندازه‌گذاری سهمیه}\label{ex:nsac}
\textbf{مسئله.} یک شرکت لجستیک در برش خود ۴۰ ارلانگ نشست هم‌زمان انتظار دارد و ۱٪ رد را می‌پذیرد. چه سهمیهٔ \acro{NSAC} باید پیکربندی
شود و اگر بار آن ۱۲٫۵٪ رشد کند احتمال رد چگونه تغییر می‌کند؟

\textbf{حل.} $c = 53$ (\cref{sec:nsac}). در $A = 45$ ارلانگ با $c = 50$ احتمال رد ۵٫۴٪ می‌شد؛ با $c = 53$ برابر ۳٫۱٪ است -- همچنان بسیار
بالاتر از ۱٪؛ با رشد مستأجر سهمیه (و قیمت) باید دوباره مذاکره شود.
\end{examplebox}

% ---------------------------------------------------------------------
\section{مقایسهٔ فناوری‌ها}
\begin{table}[htbp]
\centering\small
\caption{روش‌های متمایزکردن خدمت در شبکه‌های موبایل.}\label{tab:cmp}
\rowcolors{2}{acnmist}{white}
\begin{tabularx}{\linewidth}{@{}l Y Y Y@{}}
\toprule
\rowcolor{white} & کلاس‌های کیفیت خدمات (نسل ۴/۵) & هستهٔ اختصاصی (\acro{DECOR}، نسل ۴) & برش‌بندی شبکه (نسل ۵) \\
\midrule
دانه‌بندی     & برای هر جریان/حامل          & برای هر \acro{UE} (انتخاب هسته) & برای هر نشست \acro{PDU}، سرتاسری \\
دامنه         & رفتار با بسته‌ها            & \acro{NF}های هسته              & \acro{RAN}، انتقال، هسته، مدیریت \\
جداسازی       & ندارد (اولویت‌ها)           & نمونه‌های هسته                  & منطقی تا فیزیکی \\
استفادهٔ هر \acro{UE} & جریان‌های بسیار     & یک هستهٔ اختصاصی              & تا ۸ برش \\
دید مشتری     & عرضه نمی‌شود                & عرضه نمی‌شود                    & محصول با توافق‌نامه، مدیریت مستأجر \\
\bottomrule
\end{tabularx}
\end{table}

\paragraph{برش‌بندی و کیفیت خدمات} برش‌ها و جریان‌های کیفیت خدمات (درس ۲۲) مکمل‌اند: برش تعیین می‌کند \emph{کدام} شبکه و منابع به یک نشست
خدمت می‌دهند؛ جریان‌های کیفیت خدمات تعیین می‌کنند با بسته‌ها \emph{درون} آن چگونه رفتار شود. یک برش \acro{URLLC} همچنان از جریان‌های کیفیت
خدمات با \acro{5QI}های مناسب استفاده می‌کند.

% ---------------------------------------------------------------------
\section{نکته‌های کلیدی}
\begin{takeawaybox}
\begin{itemize}
  \item برش یک شبکهٔ منطقی سرتاسری است که با \lr{S-NSSAI} = \acro{SST} (۸ بیت) + \acro{SD} اختیاری (۲۴ بیت) شناسایی می‌شود.
  \item مجموعه‌های \acro{NSSAI}: پیکربندی‌شده، درخواستی، مجاز (حداکثر ۸)، ردشده؛ \acro{NSSF} برش‌ها و مجموعه‌های \acro{AMF} را انتخاب می‌کند.
  \item برش‌ها \acro{NF}های مشترک و اختصاصی را ترکیب می‌کنند؛ برای تضمین سرتاسری \acro{RAN} و انتقال نیز باید برش‌بندی شوند.
  \item جداسازی بهرهٔ چندگانه‌سازی را هزینه می‌کند: $G = \sum(\mu_i + z\sigma_i)/(\sum\mu_i + z\sqrt{\sum\sigma_i^2})$.
  \item بخش‌های \acro{RAN} $R_i^{\min} = \lfloor\rho_i N\rfloor r_i$ می‌دهند؛ سهمیه‌های \acro{NSAC} با ارلانگ \lr{B} اندازه‌گذاری می‌شوند.
  \item \acro{NSSAA} و \acro{NSAC} به مستأجران اجازهٔ کنترل دسترسی و به اپراتورها اجازهٔ اعمال ظرفیت قراردادی می‌دهند.
\end{itemize}
\end{takeawaybox}

% ---------------------------------------------------------------------
\section{تمرین‌ها}

\begin{exercisebox}{1}{برش یا کیفیت خدمات؟}\label{exr:slqos}
برای هر نیازمندی بگویید یک برش جداگانه یا یک جریان کیفیت خدمات درون برشی موجود ابزار بهتری است: (الف) اولویت‌دادن به تماس‌های ویدئویی
نسبت به بارگیری‌ها برای مصرف‌کنندگان؛ (ب) شبکهٔ پلیسی که باید هنگام ازدحام شبکهٔ عمومی کار کند؛ (ج) کارخانه‌ای که داده‌هایش نباید
محل را ترک کند.
\end{exercisebox}

\begin{exercisebox}{1}{سطوح جداسازی}\label{exr:iso}
برای جداسازی منطقی، مجازی و فیزیکی در هسته و در \acro{RAN} نمونه بیاورید.
\end{exercisebox}

\begin{exercisebox}{2}{برش‌های انفجاری}\label{exr:gain}
محاسبهٔ \cref{sec:mux} را با $\sigma = 60$ تکرار کنید. بهره چقدر است و برای برش‌هایی که ترافیک انفجاری حمل می‌کنند چه معنایی دارد؟
\end{exercisebox}

\begin{exercisebox}{2}{سهم \lr{RAN}}\label{exr:prb}
به یک برش \acro{URLLC} ۱۵٪ از ۲۷۳ بلوک منبع با \acro{MCS} مقاوم ۰٫۲ مگابیت بر ثانیه برای هر بلوک تضمین شده است. نرخ تضمین‌شدهٔ آن را محاسبه
کنید. چند جریان کنترلی ۶۴ کیلوبیتی جا می‌شود و اگر سهم «کمینه» باشد نه «اختصاصی»، بر سر بلوک‌های استفاده‌نشده چه می‌آید؟
\end{exercisebox}

\begin{exercisebox}{3}{طراحی سهمیه}\label{exr:nsac}
نشان دهید کوچک‌ترین سهمیه با $B \le 1\,\%$ در $A = 40$ ارلانگ برابر $c = 53$ است (از بازگشت $B(A,k) = A B(A,k-1)/(k + A B(A,k-1))$ استفاده
کنید). سپس سهمیهٔ لازم برای $A = 80$ ارلانگ را تخمین بزنید و نسبت $c/A$ را مقایسه کنید؛ روند را توضیح دهید.
\end{exercisebox}

\begin{exercisebox}{3}{طراحی یک برش}\label{exr:design}
یک بیمارستان برشی برای ۲۰۰ دستگاه متصل (پایشگرها، تصویربرداری) با تأخیر ۱۰ میلی‌ثانیه، دسترس‌پذیری ۹۹٫۹۹٪ درون پردیس و داده‌هایی که باید در
محل بمانند می‌خواهد. \lr{S-NSSAI}، مؤلفه‌های مشترک و اختصاصی، جای \acro{UPF}، اقدامات \acro{RAN} و قابلیت‌های پذیرش/مجوزدهی را پیشنهاد کنید.
\end{exercisebox}

\subsection*{پاسخ‌ها}
\addcontentsline{toc}{subsection}{پاسخ‌ها}

\begin{solutionbox}{\ref{exr:slqos}}
(الف) جریان‌های کیفیت خدمات (\acro{5QI}های متفاوت) در برش مصرفی. (ب) برشی با منابع \acro{RAN} رزروشده، اولویت و پیش‌دستی (و احتمالاً \acro{NF}های
اختصاصی هسته) -- نیازمندی، جداسازی از بار عمومی است. (ج) برشی با \acro{UPF} درون‌محلی و خروج محلی (احتمالاً \acro{NSSAA})، چون نیازمندی دربارهٔ
محل خاتمهٔ صفحهٔ کاربر است، نه اولویت بسته.
\end{solutionbox}

\begin{solutionbox}{\ref{exr:iso}}
هسته: منطقی -- \lr{SMF/UPF} مشترک با سیاست‌های هر برش؛ مجازی -- نمونه‌های اختصاصی \lr{SMF/UPF} روی کارسازهای مشترک؛ فیزیکی -- کارسازها یا سایت
اختصاصی. \acro{RAN}: منطقی -- اولویت هر برش در زمان‌بند مشترک؛ مجازی -- سهم‌های \acro{PRB} رزروشده/اختصاصی یا یک بخش پهنای باند جداگانه؛ فیزیکی --
یک حامل اختصاصی یا سلول‌های کوچک اختصاصی.
\end{solutionbox}

\begin{solutionbox}{\ref{exr:gain}}
اختصاصی $3\times(100 + 2.326\times60) = 718.7$؛ مشترک $300 + 2.326\sqrt3\times60 = 541.7$؛ $G = 1.33$. هرچه ترافیک انفجاری‌تر باشد، جداسازی ظرفیت
بیشتری هدر می‌دهد: برش‌های بهترین تلاش انفجاری باید به اشتراک بگذارند؛ جداسازی باید برای برش‌هایی نگه داشته شود که توافق‌نامه‌شان آن را لازم
می‌داند.
\end{solutionbox}

\begin{solutionbox}{\ref{exr:prb}}
$\lfloor0.15\times273\rfloor = 40$ بلوک (۴۱ با گردکردن)، $R^{\min} = 40\times0.2 = 8$ مگابیت بر ثانیه: ۱۲۵ جریان ۶۴ کیلوبیتی (با حاشیه کمتر). با سهم
\emph{کمینه}، بلوک‌های استفاده‌نشده به برش‌های دیگر قرض داده و هنگام ظهور ترافیک \acro{URLLC} پس گرفته می‌شوند (احتمالاً با تأخیر زمان‌بندی کوتاه)؛
با سهم \emph{اختصاصی} بیکار می‌مانند و قوی‌ترین تضمین را به بهای ظرفیت می‌دهند.
\end{solutionbox}

\begin{solutionbox}{\ref{exr:nsac}}
بازگشت $B(40,52) = 1.1\,\%$ و $B(40,53) = 0.82\,\%$ می‌دهد، پس $c = 53$ ($c/A = 1.33$). برای $A = 80$ ارلانگ بازگشت $c = 96$ می‌دهد ($c/A = 1.2$):
مخزن‌های بزرگ‌تر به حاشیهٔ نسبی کمتری نیاز دارند -- همان صرفهٔ مقیاس در اندازه‌گذاری ترانک (درس ۲) و بهرهٔ چندگانه‌سازی \cref{sec:mux}.
\end{solutionbox}

\begin{solutionbox}{\ref{exr:design}}
\lr{S-NSSAI}: \lr{SST 2} (\acro{URLLC}) یا یک \acro{SST} ویژهٔ اپراتور، با یک \acro{SD} که بیمارستان را مشخص می‌کند. مشترک: \acro{AMF}، \acro{NSSF}،
\acro{UDM}. اختصاصی: \acro{SMF} و یک \acro{UPF} در پردیس بیمارستان (خروج محلی، داده در محل)، احتمالاً \acro{UPF}های افزونه برای ۹۹٫۹۹٪. \acro{RAN}:
سلول‌های کوچک داخلی با سهم \acro{PRB} کمینه یا اختصاصی و اولویت برش؛ تأخیر با عددشناسی ۳۰ کیلوهرتز یا شکاف‌های کوچک. دسترسی: \acro{NSSAA} با کارساز
\acro{AAA} بیمارستان؛ سهمیهٔ \acro{NSAC} حدود ۲۰۰ \acro{UE}؛ مسیر انتقال با پهنای باند تضمین‌شده. پایش تأخیر و دسترس‌پذیری در برابر توافق‌نامه.
\end{solutionbox}

% ---------------------------------------------------------------------
\section{مطالعهٔ بیشتر}
برش‌بندی در سامانهٔ نسل پنجم در \lr{TS 23.501} \cite{ts23501} (بند ۵٫۱۵) با رویه‌ها در \lr{TS 23.502} \cite{ts23502} مشخص شده است؛ کدگذاری
\lr{S-NSSAI} در \lr{TS 23.003} \cite{ts23003}؛ مفاهیم مدیریت در \lr{TS 28.530} \cite{ts28530}؛ نیازمندی‌های خدمت در \lr{TS 22.261} \cite{ts22261}.
فوکاس و همکاران \cite{foukas2017slicing} معماری‌ها و پرسش‌های پژوهشی برش‌بندی را مرور می‌کنند؛ رومر و همکاران \cite{rommer2020core} برش‌بندی را در هستهٔ
نسل پنجم پوشش می‌دهند.

\begin{latin}
\printbibliography[heading=subbibliography,title={References}]
\end{latin}
\end{document}
